\documentclass[a4paper]{arrowhead}

\usepackage[yyyymmdd]{datetime}
\usepackage{etoolbox}
\usepackage[utf8]{inputenc}
\usepackage[T1]{fontenc}
\usepackage{multirow}
\usepackage{longtable}
\usepackage{tikz}
\usetikzlibrary{positioning,shapes.multipart,arrows.meta}

\renewcommand{\dateseparator}{-}
\setlength{\emergencystretch}{3em}
\lstset{language=HTTP,breakatwhitespace=false,columns=fullflexible}

%% Special references
\newcommand{\mref}[2][]{{\textcolor{ArrowheadPurple}{\hyperref[sec:model:#2]{\ifstrempty{#1}{#2}{#1}}}}}
\newcommand{\pref}[1]{{\textcolor{ArrowheadGrey}{\hyperref[sec:model:primitives:#1]{#1}}}}
%%

\begin{document}
\sloppy

%% Arrowhead Document Properties
\ArrowheadTitle{LocalCloudManager}
\ArrowheadType{System Design Description}
\ArrowheadTypeShort{SysDD}
\ArrowheadVersion{5.2.2}
\ArrowheadDate{2026-11-09}
\ArrowheadAuthor{Jerker Delsing}
\ArrowheadStatus{RELEASE CANDIDATE}
\ArrowheadContact{jerker.delsing@ltu.se}
\ArrowheadFooter{\href{www.arrowhead.eu}{www.arrowhead.eu}}
\ArrowheadSetup
%%

%% Front Page
\begin{center}
  \vspace*{1cm}
  \huge{\arrowtitle}

  \vspace*{0.2cm}
  \LARGE{\arrowtype}
  \vspace*{1cm}

  \vspace*{\fill}

  % Front Page Image

  \vspace*{1cm}
  \vspace*{\fill}

  % Front Page Abstract
  \begin{abstract}
    This document describes the design and implementation of the LocalCloudManager support system of Arrowhead 5.2.2 (\texttt{ah5-LocalCloudManager-java-spring}).
  \end{abstract}

  \vspace*{1cm}
\end{center}

\newpage
%%

%% Table of Contents
\tableofcontents
\newpage
%%

\section{Overview}
\label{sec:overview}

This document describes the LocalCloudManager system implementation, which realizes the LocalCloudManager SysD \cite{lcm-sysd}. LocalCloudManager is the operator console of the local cloud: it shows the state of the cloud and carries out the management actions an operator starts, with its own Arrowhead identity.
In Section \ref{Implementation}, we describe implementation details of the system.

\newpage

\section{Implementation}
\label{Implementation}

This implementation is based on the SysD document LocalCloudManager \cite{lcm-sysd} and is located in \texttt{ah5-LocalCloudManager-java-spring} of the primary Java source tree (\texttt{Eclipse-Arrowhead-java}). It is a Java-Spring support system built from two Maven projects: the console and Arrowhead system (artifact \texttt{arrowhead-\allowbreak{}localcloudmanager}, \texttt{localcloudmanager.\allowbreak{}jar}) and the Docker proxy (\texttt{docker-\allowbreak{}proxy/}, artifact \texttt{arrowhead-\allowbreak{}localcloudmanager-\allowbreak{}docker-\allowbreak{}proxy}).

\subsection{Implementation language and tools}
\begin{itemize}
\item Java 21, Spring Boot 3.4 (Web, Security, WebFlux client, Quartz), Log4j2, BouncyCastle (scrypt).
\item Eclipse Arrowhead common library as a private vendored source copy (\texttt{vendor/}, the variant without JPA and log database): system information, ServiceRegistry registration, outsourced authentication and management filters, service collector, sysMonitor publisher, security profile guard.
\item Storage: none (no database).
\item Docker proxy: minimal Spring Boot Web application, Unix-socket HTTP client of the Docker Engine API.
\end{itemize}

\subsection{Functional properties implementation}

The following components implement the system (package \texttt{eu.\allowbreak{}arrowhead.\allowbreak{}localcloudmanager}):

\noindent\begin{longtable}{| >{\raggedright\arraybackslash}p{5.6cm} | >{\raggedright\arraybackslash}p{9.6cm} |} \hline
\rowcolor{gray!33} Component & Responsibility \\ \hline \endhead
\texttt{LocalCloud\allowbreak{}Manager\allowbreak{}Main}, \texttt{LocalCloud\allowbreak{}Manager\allowbreak{}SystemInfo} & Application; system \texttt{LocalCloudManager} (metadata \texttt{role}, \texttt{managedProject}, \texttt{consumerInterfaceTemplates}, X.509 public key) and the services \texttt{local\allowbreak{}Cloud\allowbreak{}Management} and \texttt{general\allowbreak{}Management}. \\ \hline
\texttt{init/\allowbreak{}LocalCloud\allowbreak{}Manager\allowbreak{}Application\allowbreak{}Init\allowbreak{}Listener} & Registration by the common listener; registration state for \texttt{/health} and the overview. \\ \hline
\texttt{console/\allowbreak{}ConsoleConnectorConfig} & Second Tomcat connector (plain HTTP, \texttt{console.\allowbreak{}port}, \texttt{console.\allowbreak{}address}). \\ \hline
\texttt{console/\allowbreak{}ConsoleFilter} & Separates the connectors: on 8490 the static UI (\texttt{ui/\allowbreak{}public}, packaged as \texttt{console/}), \texttt{/health}, proxy-secret, Host/Origin and password-configured checks, marks the Arrowhead filters as done; on 8491 only \texttt{/api}, \texttt{general\allowbreak{}Management} and the API documentation. \\ \hline
\texttt{console/\allowbreak{}ConsoleSecurityConfig} & Spring Security: console chain (session cookie as authentication \texttt{ROLE\_\allowbreak{}OPERATOR}, \texttt{/api} except \texttt{/api/\allowbreak{}session} needs it, headers, stateless) and Arrowhead chain (HTTPS only). Replaces the common \texttt{Default\allowbreak{}Security\allowbreak{}Config}. \\ \hline
\texttt{console/\allowbreak{}ConsoleAuth} & Console login: scrypt password hash, sessions, per-client lock-out, Host/Origin check, proxy secret. \\ \hline
\texttt{api/\allowbreak{}http/\allowbreak{}ConsoleApiController} & All \texttt{/api} routes and \texttt{/health}, matched in the order of the native console. \\ \hline
\texttt{api/\allowbreak{}http/\allowbreak{}GeneralManagementAPI} & \texttt{general\allowbreak{}Management} \texttt{get-\allowbreak{}config}. \\ \hline
\texttt{security/\allowbreak{}Manager\allowbreak{}AccessFilter}, \texttt{Manager\allowbreak{}Authorization\allowbreak{}Service} & On 8491, after authentication: Sysop, \texttt{lcm.\allowbreak{}api.\allowbreak{}whitelist} or ConsumerAuthorization grant of the \texttt{local\allowbreak{}Cloud\allowbreak{}Management} operation (cached 30 s, fail closed). \\ \hline
\texttt{client/\allowbreak{}Orchestrated\allowbreak{}Provider\allowbreak{}Gateway} & Providers selected by the common service collector (orchestration on behalf of LocalCloudManager); HTTPS with the identity token; status and body passed through for LocalCloudConfigurator. \\ \hline
\texttt{client/\allowbreak{}TlsContextFactory} & Client TLS context: system certificate, cloud truststore, CRL of the cloud CA (rebuilt when the CRL file changes). \\ \hline
\texttt{client/\allowbreak{}DockerGateway} & HTTP client of the Docker proxy. \\ \hline
\texttt{service/\allowbreak{}ManagerService} & Views and actions: overview, registry management, DataManager, consumers, inter-cloud, sysMonitor, plant descriptions and SoS, orchestration rules and history, policies, logs, lifecycle, floorplans, local-cloud configuration. \\ \hline
\texttt{service/\allowbreak{}Inventory} & Inventory names, orchestration catalog, rule validation, registry management, messaging overview, SenML resolution, Docker log frames. \\ \hline
\texttt{service/\allowbreak{}PolicyEditor}, \texttt{RuleEditor}, \texttt{SosModel} & Policy grant/replace/revoke with snapshot checks; rule priority (native or create-and-remove); SoS validation, inclusion and cycle checks. \\ \hline
\texttt{service/\allowbreak{}CrlStatus}, \texttt{Docker\allowbreak{}Health\allowbreak{}Indicator} & Remaining CRL validity; sysMonitor health (degraded while Docker is unreachable). \\ \hline
\texttt{docker-\allowbreak{}proxy/\allowbreak{}Allowlist}, \texttt{Proxy\allowbreak{}Controller}, \texttt{Docker\allowbreak{}Socket} & Docker proxy: decision per request, forwarding of allowed requests to \texttt{/\allowbreak{}var/\allowbreak{}run/\allowbreak{}docker.\allowbreak{}sock}. \\ \hline
\end{longtable}

A console request passes the \texttt{ConsoleFilter} (proxy secret for \texttt{/api}, Host/Origin, password configured), the Spring Security console chain (session) and, for changes, the confirmation check before the route handler calls the providing systems. A request on the Arrowhead connector passes the common authentication filter (client certificate and identity token), the management filter (\texttt{general\allowbreak{}Management}) and the \texttt{ManagerAccessFilter}.

\subsubsection{Data handled and stored}

None durable. Console sessions, the identity token, orchestration results (300 s) and the last 60 overview samples are kept in memory. SoS compositions, policies, rules and configurator designs are stored by their providing systems.

The interfaces exchange JSON objects built at run time; their structure is specified in the IDD documents.

\subsection{Non functional properties implementation}

\subsubsection{Security}

The console process has no Docker socket; Docker is reached only through the separate proxy on the internal \texttt{control} network, which forwards only the listed requests (SysD, Section Security). The console connector is published on the host's loopback address only. The security profile guard refuses insecure settings at start-up. Both containers run with a read-only root filesystem and without capabilities; the console as UID 10001.

\subsubsection{Power management}

None; the system is a server process without power states. The \texttt{sysMonitor} \texttt{power} observation always reports line supply.

\subsubsection{Internal monitoring}

Logs go to the console and a rolling file (Log4j2, no database appender). The common publisher sends the \texttt{sysMonitor} observations every 30 s (see the sysMonitor IDD); the echo is degraded while the Docker proxy cannot be reached.

\subsubsection{Configuration}

Spring properties (listed in the SysD), with secrets from \texttt{configtree:\allowbreak{}/\allowbreak{}run/\allowbreak{}secrets/}.

\subsection{Deployment}

The local \texttt{compose.yaml} (project \texttt{ah5-\allowbreak{}localcloudmanager}) builds and runs two containers from this folder:

\begin{itemize}
\item \texttt{cloud-\allowbreak{}monitor-\allowbreak{}api} (image \texttt{arrowhead-\allowbreak{}localcloudmanager-\allowbreak{}standalone:\allowbreak{}5.2.2}): ports 8490 (published on \texttt{127.0.0.1}) and 8491, on the external cloud network and the internal \texttt{control} network; key store, truststore and CRL mounted read-only; secrets store password, identity password and console password hash.
\item \texttt{docker-\allowbreak{}proxy} (image \texttt{arrowhead-\allowbreak{}localcloudmanager-\allowbreak{}docker-\allowbreak{}proxy-\allowbreak{}standalone:\allowbreak{}5.2.2}): port 2375 on the \texttt{control} network only, Docker socket mounted read-only, \texttt{COMPOSE\_\allowbreak{}PROJECT\_\allowbreak{}NAME} and \texttt{MANAGED\_\allowbreak{}SERVICES}.
\end{itemize}

The earlier combined test cloud (\texttt{Arrowhead-test-cloud-java}) is a separate deployment consumer and still runs the earlier packaging until it is redeployed.

\subsection{Local modifications}
\label{sec:local}

Notes on this implementation in the workspace:
\begin{itemize}
  \item Workspace system; operation names are camelCase (see the SD).
  \item \texttt{session}, \texttt{login} and \texttt{logout} are console-only and no longer registered; console endpoints for policies, rule priority, SoS, local-cloud configuration and chiplet floorplans are not part of the registered operation list.
  \item \texttt{general\allowbreak{}Management} has no \texttt{get-\allowbreak{}log} (no log database).
\end{itemize}

\subsection{Verification}

\texttt{mvn clean verify} runs the JUnit tests of both projects: a full-application MockMvc test of every \texttt{/api} route on both connectors (console login flow, Host/Origin refusal, confirmation header, pass-through of LocalCloudConfigurator, identity and grant checks on 8491, \texttt{general\allowbreak{}Management}); a real-Tomcat test that the console connector listens on the loopback address only and serves \texttt{/api} only after login; console authentication, CRL status, inventory, policy, rule-priority and SoS tests (ports of the native Node tests); the security profile test; and the Docker proxy allowlist (MockMvc with a stub socket). The browser models in \texttt{ui/\allowbreak{}test} are Node tests of the UI code.

\newpage

\section{Services}
\label{sec:services}

The implementation of the services is based on the SD and IDD documents referenced below.

\noindent\begin{longtable}{| p{5.4cm} | p{2.2cm} | p{2.2cm} | p{4.8cm} |} \hline
\rowcolor{gray!33}\textbf{Services produced} & \textbf{SysD ref} & \textbf{SD ref} & \textbf{IDD ref} \\ \hline \endhead
\texttt{local\allowbreak{}Cloud\allowbreak{}Management} & \cite{lcm-sysd} & \cite{lcm-localCloudManagement-sd} & \cite{lcm-localCloudManagement-http-idd} \\ \hline
\texttt{sysMonitor} & \cite{lcm-sysd} & \cite{lcm-sysMonitor-sd} & \cite{lcm-sysMonitor-amqp-idd} \\ \hline
\texttt{general\allowbreak{}Management} & \cite{lcm-sysd} & -- & -- \\ \hline
&&& \\ \hline
\rowcolor{gray!33}\textbf{Services consumed} & \textbf{SysD ref} & \textbf{SD ref} & \textbf{IDD ref} \\ \hline
serviceDiscovery; service\allowbreak{}Registry\allowbreak{}Management & \cite{core-sr-sysd} & \cite{core-sr-serviceDiscovery-sd} & -- \\ \hline
identity & \cite{core-auth-sysd} & \cite{core-auth-identity-sd} & -- \\ \hline
service\allowbreak{}Orchestration; service\allowbreak{}Orchestration\allowbreak{}History\allowbreak{}Management & \cite{core-dso-sysd} & \cite{core-dso-serviceOrchestration-sd} & -- \\ \hline
service\allowbreak{}Orchestration\allowbreak{}Simple\allowbreak{}Store\allowbreak{}Management & \cite{core-ssso-sysd} & \cite{core-ssso-serviceOrchestrationSimpleStoreManagement-sd} & -- \\ \hline
authorization\allowbreak{}Management (query, grant, revoke policies) & \cite{core-ca-sysd} & \cite{core-ca-authorizationManagement-sd} & -- \\ \hline
dataHistorian & \cite{dm-sysd} & \cite{dm-dataHistorian-sd} & -- \\ \hline
sysMonitorHistory; sys\allowbreak{}Monitor\allowbreak{}Management & \cite{smo-sysd} & \cite{smo-sysMonitorHistory-sd} & -- \\ \hline
plant\allowbreak{}Description\allowbreak{}Management & \cite{pde-sysd} & \cite{pde-plantDescriptionManagement-sd} & -- \\ \hline
gatekeeper\allowbreak{}Management; gate\allowbreak{}Tunnel\allowbreak{}Management & Gatekeeper, GateTunnel & -- & -- \\ \hline
Docker Engine subset & Docker proxy & -- & -- \\ \hline
local\allowbreak{}Cloud\allowbreak{}Configuration & \cite{lcc-sysd} & \cite{lcc-localCloudConfiguration-sd} & \cite{lcc-localCloudConfiguration-http-idd} \\ \hline
floorplan\allowbreak{}Discovery; floorplan\allowbreak{}Management & Chiplet\allowbreak{}Floorplan\allowbreak{}Generator & -- & -- \\ \hline
\end{longtable}


\bibliographystyle{IEEEtran}
\bibliography{bibliography}

\newpage

\section{Revision History}
\subsection{Amendments}

\noindent\begin{tabularx}{\textwidth}{| p{1cm} | p{3cm} | p{2cm} | X | p{4cm} |} \hline
\rowcolor{gray!33} No. & Date & Version & Subject of Amendments & Author \\ \hline

1 & 2026-10-22 & \arrowversion & Document created from the Arrowhead LaTeX templates for the 5.2.2 Java Spring implementation, including the local extensions of this workspace. & Jerker Delsing \\ \hline
2 & 2026-10-30 & \arrowversion & Checked against the primary Java source tree (Eclipse-Arrowhead-java): location, new console modules, standalone deployment, consumed services and tests updated. & Claude Opus 5.5 (draft) \\ \hline
3 & 2026-11-09 & \arrowversion & Java-Spring implementation: components, connectors (8490 console, 8491 Arrowhead), Spring Security console login, Docker proxy project, orchestrated providers including LocalCloudConfigurator, deployment and JUnit/MockMvc verification. Node.js/nginx descriptions removed. & Claude Opus 5.5 (draft) \\ \hline
\end{tabularx}

\subsection{Quality Assurance}

\noindent\begin{tabularx}{\textwidth}{| p{1cm} | p{3cm} | p{2cm} | X |} \hline
\rowcolor{gray!33} No. & Date & Version & Approved by \\ \hline

1 & & \arrowversion  &  \\ \hline

\end{tabularx}

\end{document}
